\documentclass{article}
\usepackage{pathfinder-note}

\title{Cost-Aware Endogenous Advice in The Energy Society}

\begin{document}
\maketitle

\section{Pair and connexion}

Q surveys learning-augmented algorithms, including prediction cost, feedback, composition, endogenous error, and semantic predictors. P introduces a survival economy in which inference consumes energy, incentives alter collective behavior, and agent recommendations affect coordination and job selection. The strongest connexion is therefore an empirical study of cost-adjusted, endogenous advice under strategic feedback. It is not a demonstrated transfer of a theorem from Q into P \ledger{1,7,13,18}.

\section{Supported findings}

The abstracts support one focused research hypothesis: competitive and cooperative incentives may shift the consistency--robustness frontier of advice after accounting for inference cost. A suitable experiment would compare a fixed advice-free fallback with accurate and corrupted exogenous advice and with endogenous advice generated under randomized incentive regimes. Primary outcomes should include group survival and net energy or completed-job welfare after charging all recommendation-token costs \ledger{3,10,16,21}.

The primary analysis should be intention-to-treat: randomize incentive regime, recommendation availability, and corruption ex ante, then estimate their effects and interactions relative to the fallback. Trace replay can provide mechanism evidence, but may disrupt history-dependent feasibility. Comparisons conditioned on realized prediction error are descriptive because error is itself affected by incentives and feedback \ledger{14,19}.

\section{Objections and limits}

The evidence is thin: only abstracts are available, and no experimental result for this connexion has been established. Q supplies neither a theorem statement nor a concrete prediction interface, error measure, or fallback in the provided text. P supplies neither a formal welfare guarantee nor evidence of an offline comparator or the intervention hooks required by the proposal. Consequently, present-tense claims that Q's framework has been implemented or that its guarantees are preserved are unsupported \ledger{2,9,11}.

There is also a construct-validity problem. An action recommendation is not automatically a prediction. Recommendations should be represented as forecast--action tuples, such as a job choice paired with predicted success probability or net-energy return. Forecast quality should be scored separately from policy value. Because outcomes may be observed only for selected jobs, any calibration claim must address selective labels through randomized audits, sufficiently complete outcome logging, or an explicit restriction to on-policy accuracy \ledger{15,20}.

\section{Unresolved issues}

The study still requires Q's full prediction interface and sufficient conditions; a meaningful prediction-independent fallback in P; complete inference-cost accounting; feasible semantic corruption interventions; seeded paired trials and uncertainty intervals; and a defensible welfare comparator. If these conditions cannot be instantiated, the correct conclusion is that the proposed connexion yields no result, rather than that empirical averages constitute a formal guarantee \ledger{4,16,21}.

\section{Smallest next step}

Inspect the full Q paper and P implementation together. Write down one forecast--action interface, one deterministic advice-free fallback, and one fully costed randomized intervention that the simulator can execute. Only after this feasibility check should the exogenous-versus-endogenous frontier experiment be preregistered. A positive result would show net gains under accurate advice with bounded loss under corruption; a negative result would show an intention-to-treat frontier collapse under competitive endogenous advice after controlling token cost and corruption assignment \ledger{19,21}.

\end{document}